\section{System sygnalizacji pożarowej SSP/SAP}
\subsection{Rodzina urządzeń i rozmieszczenie}
CSP-0 to SATEL ACSP-402 z fabrycznym zasilaczem APS-318, umieszczona w stale obsadzonym SOC. Zastosowano \FireDetectorCount{} czujek: dziesięć DMP-400 dymu i ciepła oraz jedną DCP-400 ciepła A1R w WC, \RopCount{} ROP-400/PL i \FireSounderCount{} SPP-400. Wszystkie elementy adresowalne należą do rodziny ACSP, mają obustronne izolatory i zasilanie z pętli. Gniazda DB-400 wchodzą w zespoły czujek. Nie ma sieci między centralami ani interfejsu światłowodowego.\cite{S-FIRE,S-DMP,S-ROP,S-SND}

Dla płaskiego stropu na 3 m przyjęto własny promień kontrolny 5 m dla detekcji dymu i 4 m dla ciepła w WC. Nie są to parametry producenta ani cytat z pełnego tekstu normy rozmieszczenia. Każdy narożnik pomieszczenia i siatka kontrolna co 0,5 m leżą w takim promieniu od przypisanej czujki. W holu 16 $\times$ 10 m cztery czujki wynikają z podziału na obszary 8 $\times$ 5 m; najdalszy narożnik jest oddalony o 4,72 m. W przejściu 8 $\times$ 10 m wystarczają dwa takie obszary. Dla pozostałych pomieszczeń najdalsze narożniki są w odległości 3,61--4,24 m; czujka ciepła w WC ma odległość maksymalną 3,61 m.

DMP w IT znajduje się nad wolnym przejściem przed szafą, a w magazynie nad środkowym przejściem między regałami. Nie ma pustek wymagających osobnych czujek. Przyjęto brak nawiewów przy punktach; zmiana warunków wentylacji lub dodanie przegród podsufitowych jest zmianą modelu, a nie korektą rysunku. Czujkę ciepła w WC wybrano wobec możliwej pary wodnej; przyjęto brak prysznica i procesu wytwarzającego gorące powietrze. Dozór cieplny może reagować później od dymowego.

\FirePointTable

\subsection{Pętla, adresy i sygnalizacja}
L1 jest pętlą zamkniętą CSP-0:L1A--elementy według adresów 1--22--CSP-0:L1B; dokładna kolejność kabli FL1-01--FL1-23 jest w zestawieniu połączeń. Nie są to odgałęzienia siatki orientacyjnej. Zajęto 22 z 128 adresów, wszystkie z izolatorami. Druga pętla pozostaje wyłączona, bez zakupionych punktów rezerwowych. Wyjście i powrót prowadzi się na odrębnych uchwytach; powrót w SOC prowadzi stroną północną, a pierwsze odejście w stronę holu. Pojedyncza przerwa pozwala na zasilanie z drugiej strony; zwarcie izoluje fragment między sąsiednimi izolatorami. Nie deklaruje się odporności na dwa dowolne uszkodzenia ani utratę samej CSP.\cite{S-FIRE,S-FIRE-USER}

ROP-y umieszczono przy E01, E02 i E03 na wysokości 1,4 m. Kontrola najkrótszej drogi przez rzeczywiste otwory drzwiowe daje do około 19,5 m wobec własnego celu 30 m; odległość nie jest mierzona przez ścianę. Po jednym sygnalizatorze przyjęto w każdym zamkniętym pomieszczeniu, a dwa w holu. Przyjęto tło do 55 dBA i cel co najmniej 65 dBA lub 5 dB ponad tłem w punktach użytkowania przy zamkniętych drzwiach. Nie jest to wynik pomiaru. Model ma sygnalizację akustyczną; nie deklaruje pokrycia optycznego EN 54-23. Osobie, która nie odbiera sygnału akustycznego, pomaga wyznaczony pracownik. Odbiór obejmuje ocenę dostępności i próbę tej procedury.

\subsection{Macierz przyczyna--skutek i wyjścia}
Przyjęto jednostopniowe alarmowanie: dowolna czujka albo ROP uruchamia alarm ewakuacyjny bez oczekiwania na AI i bez czasu na potwierdzenie przez SOC. Nie projektuje się oddymiania, klap, wind, gaszenia ani wspólnej wentylacji mechanicznej. Nie dopisuje się wyjść do instalacji, których nie ma w modelu.

\begin{tabularx}{\linewidth}{@{}P{29mm}P{38mm}Y@{}}
\toprule Wyzwalacz / warunek & Rzeczywisty kanał & Skutek, nadzór i przywrócenie\\\midrule
F01--F11, alarm bez zwłoki & L1: adresy 3, 8, 10, 14, 16, 18, 20, 22 (SF-01--08); OUT1. & Wszystkie sygnalizatory; ekran CSP; informacja do Z11 SSWiN. Odpytanie pętli i izolatory, linia INFO-FIRE 2EOL. Reset lokalny na poziomie 2 po sprawdzeniu przyczyny i zgodzie dyżurnego.\\
ROP-01--03, naciśnięcie & Te same osiem adresów i OUT1. & Niezwłoczny alarm. ROP najpierw resetowany kluczem, następnie CSP po potwierdzeniu bezpieczeństwa.\\
Przerwa, zwarcie, brak elementu, uszkodzenie zasilania/baterii & OUT8, przekaźnik wzbudzony w stanie sprawnym. & Informacja do Z12 i lokalne wskazanie CSP. Zanik zasilania CSP także powoduje stan uszkodzenia; przewód 2EOL odróżnia przerwę/zwarcie. Przywrócenie po naprawie i sprawdzeniu nadzoru.\\
Pożar i zamknięte KD & Brak elektrycznego wyjścia zwalniania. & Wewnętrzne klamki D01--D04 i okucia E01--E03 otwierają mechanicznie. Kontrola fizyczna po zdarzeniu; brak zależności od LAN.\\
Brak LAN/VMS/AI & Brak sterowania SSP z IT. & Lokalny alarm i obsługa CSP trwają. Powrót IT po naprawie; nie kasuje automatycznie alarmu pożarowego.\\\bottomrule
\end{tabularx}

Do opisanych funkcji wykorzystuje się osiem adresowalnych sygnalizatorów i dwa fizyczne styki informacyjne centrali: OUT1 i OUT8. Nie zakupiono modułów MIO, bo żadna przyjęta funkcja ich nie wymaga. Niewykorzystane fabryczne wyjścia nie są liczone jako instalowane sterowania. Linie informacyjne do SSWiN nie są torem transmisji alarmu do PSP.\cite{S-FIRE,S-FIRE-USER}

\subsection{DSO i powiadamianie straży}
Dla całego przyjętego fikcyjnego portu ograniczono jednoczesną obsadę do 150 osób. Nie przekracza to progu ponad 500 osób dla portów w §28 ust.1 pkt19 i §29 ust.1 pkt8 rozporządzenia. Model nie spełnia też pozostałych przyjętych przesłanek tych katalogów: jest niski, parterowy, bez funkcji noclegowej, medycznej, archiwalnej, garażu podziemnego ani dużego obiektu handlowego lub widowiskowego. SSP przyjęto dobrowolnie ze względu na zakres projektu; DSO nie instaluje się. Wniosek dotyczy zdefiniowanego małego portu, a nie dowolnego lotniska lub dowolnie wyciętego fragmentu dużego terminalu.\cite{L-FIRE}

Przejrzano zmiany Dz.U.2024 poz.1716 (m.in. oznakowanie i autonomiczne czujki) i Dz.U.2026 poz.1276 (oznakowanie punktów schronienia, obowiązuje od 30.09.2026); nie zmieniają wskazanych progów. Art.5 ustawy o ochronie przeciwpożarowej wiąże obowiązek połączenia do PSP z obligatoryjnym SSP; zmiana Dz.U.2026 poz.815 dotyczy art.14fa ust.3. Dla tego modelu nie przyjęto automatycznej transmisji do PSP. Dyżurny po alarmie niezwłocznie wzywa pomoc telefonicznie pod 112 i rozpoczyna ewakuację. Telefon ani ETHM nie są deklarowane jako system transmisji alarmów pożarowych. Ewentualne zwiększenie obsady całego portu lub zmiana funkcji wymaga ponownej oceny tych rozstrzygnięć.\cite{L-2024,L-2026,L-ACT,L-ACT-2026}
